% ============================================================
% CHAPTER 1: INTRODUCTION
% File: chapters/01_introduction.tex
% ============================================================

\chapter{Introduction}
\label{ch:introduction}


% ============================================================
% 1.1 CONTEXT AND MOTIVATION
% ============================================================

\section{Context and Motivation}

Industrial monitoring means tracking the state and behavior of a
physical process by analyzing data collected during or after it
runs. Industrial and robotic systems produce many types of data,
such as visual observations, acoustic signals, force measurements,
vibration, and other sensor signals
\parencite{Huo2022Multisensor,Thoduka2024Handover}.
Each data type shows a different aspect of the process and may
respond differently to changes in operation.

Multimodal learning enables the integration of these diverse data
types, particularly when each source provides complementary or
unique information. However, adding data sources does not inherently
improve predictions \parencite{Wang2020What}. The benefit depends on
the information each modality retains, the relevance of the data
to the prediction task, and the methods used for data integration.

Robotic welding was the starting point for this research. A camera
can observe the weld pool, while acoustic and machine sensor
measurements give information about other parts of the welding
process. Previous studies have explored combining visual and
acoustic data for weld assessment \parencite{P003,Hua2026Weld}.

The primary architectural challenge extends beyond welding.
Industrial monitoring systems require data representations that
capture task-relevant information, fusion methods that effectively
utilize complementary data from multiple sources, and evaluation
protocols that assess whether predictions remain reliable when some
data are missing or of low quality. Therefore, this research
investigates multimodal fusion architectures across a range of
industrial and robotic monitoring tasks.


% ============================================================
% 1.2 PROBLEM STATEMENT AND RESEARCH GAP
% ============================================================

\section{Problem Statement and Research Gap}

The added value of additional modalities cannot be determined
solely by the predictive performance of the combined model; it is
essential to compare the architecture's performance against
credible unimodal predictors on identical prediction targets
and assessment observations.

Prior work has shown that multimodal integration can improve
prediction, but it also identifies circumstances in which a jointly
trained multimodal model will perform worse than its unimodal
counterparts \parencite{Wang2020What}. In one weld penetration study,
the image-only, audio-only, and fused models all achieved 100\%
accuracy on the reported clean test set
\parencite[Tables~1--3]{Hua2026Weld}. There was therefore no demonstrable
benefit to the availability of an additional modality in terms
of clean data accuracy.

Another challenge concerns the reliability of each data type.
When an input is missing, fallback mechanisms may be employed.
If an input is present but degraded, evidence is required to
demonstrate altered reliability. While adaptive fusion can adjust
the contribution of each input, modifying these weights alone
does not ensure improved robustness \parencite{Zhang2023QMF,P008}.

The fairness of these comparisons also depends on factors like
the prediction target, how the data is aligned over time, training
methods, calibration, and whether the test data is independent
\parencite{P017,P004,Guo2017Calibration,P008}. Differences in any of
these areas can make it hard to know if performance gains are
generalizable.

This thesis addresses this gap by investigating the added predictive
value and robustness of multimodal fusion under fair, task-specific
experimental conditions. It investigates how the selection of data
representation, complementary information, and fusion method
interact, and whether observed benefits in one context are
transferable to other industrial and robotic monitoring tasks.

Chapter~2 reviews existing research to lay out the scientific
background and highlight the unresolved comparisons that define
this research gap.


% ============================================================
% 1.3 RESEARCH QUESTIONS AND OBJECTIVES
% ============================================================

\section{Research Questions and Objectives}

The main goal of this thesis is to study how to design and evaluate
process-aware multimodal AI systems for reliable industrial
monitoring and anomaly detection. It examines how using specific
data representations and different fusion methods affects prediction,
robustness, and practical usefulness.

The first version of the expos\'{e} set out this main research
question:

\begin{quote}

How can a process-aware multimodal AI architecture be designed
and evaluated for robust industrial anomaly detection using
robotic welding data as a case study?

\end{quote}

Four subquestions address the main architectural and methodological
aspects of this study.

\begin{enumerate}[label=SQ\arabic*.,leftmargin=16mm]

\item Which lightweight or pretrained Transformer-based
representation components are suitable for visual, audio, and
sensor time-series data in a multimodal industrial anomaly
detection pipeline?

\item How do unimodal baselines, simple concatenation fusion,
modular late fusion, and process-aware gated cross-modal fusion
compare in terms of predictive performance?

\item Can a gated cross-modal fusion architecture improve
robustness when one or more modalities are missing, degraded,
or unreliable?

\item How can multimodal architecture performance be evaluated
beyond accuracy, including robustness, interpretability,
modularity, and practical usability?

\end{enumerate}

\ev{Research questions and scope}
{\EvidenceRef{E002}{pp. 2--3}, \EvidenceRef{E001}{}}

These considerations result in four research objectives: evaluating
the suitability of specific data representations, investigating
various fusion strategies, testing robustness to input degradation,
and assessing the appropriateness of selected evaluation methods
for multimodal monitoring systems. Broader architectural questions
persist, despite changes in the original application scope.


% ============================================================
% 1.4 METHODOLOGY OVERVIEW AND CONTRIBUTIONS
% ============================================================

\section{Methodology Overview and Contributions}

This research integrates a critical literature review with empirical
investigations of multimodal representation and fusion architectures.
The review encompassed \ReviewWorkCount{} full texts, screened and
assessed in two stages, with emphasis on architectural approaches,
unimodal baselines, fusion methods, evaluation protocols, and
strategies for handling unreliable inputs.

Dataset selection also considered modality correspondence,
annotation suitability, assessment independence, and computational
viability. The empirical component compared modality-specific
representations and prediction models, then evaluated alternative
fusion strategies and a series of controlled follow-up experiments.

The implemented architectures use pre-trained convolutional visual
representations, statistical sensor features, and patch-based
temporal encoders. Information is combined using feature
concatenation, probability averaging, logistic stacking, and
context-conditioned weighting.

Evaluation covers predictive performance, calibration, grouped
assessment, uncertainty, and how the system behaves when inputs
are missing or degraded, depending on the dataset and protocol.
Tests on external datasets check if the chosen fusion methods
and data representations work for different prediction tasks.

The principal contributions of this thesis are:

\begin{itemize}

\item \textbf{Comparative evaluation of multimodal representations:}
This includes empirical tests of pretrained visual features,
statistical sensor representations, and patch-based temporal
encoders. The study compares training budgets, feature combinations,
and the compatibility of different representations. It assesses
each representation's performance individually and in combination.

\item \textbf{Comparative evaluation of fusion architectures:}
Feature-level, prediction-level, and process-aware fusion
architectures are implemented and evaluated against unimodal
reference models. These comparisons establish the conditional
predictive gain from information fusion and distinguish this
gain from the effects of additional fitting and calibration
to the extent permitted by the experimental controls.

\item \textbf{Assessment of multimodal robustness:}
The experiments test how fusion behaves when some data is
missing or degraded, comparing fallback methods to adaptive
input weighting. The study checks if the benefits seen with
clean data remain when inputs are degraded.

\item \textbf{Cross-dataset evaluation and methodological assessment:}
The research explores re-using fusion methods and transferring
supervised representations across different industrial and
robotic monitoring tasks. The evaluation highlights the
importance of target validity, data alignment, class-specific
performance, grouped assessment, and uncertainty when comparing
architectures across datasets.

\end{itemize}

Implementation records, literature comparison tables, and
experimental artifacts all provide clear support for these
contributions.


% ============================================================
% 1.5 SCOPE AND DELIMITATIONS
% ============================================================

\section{Scope and Delimitations}

The original research plan focused on robotic welding as the main
application for process-aware multimodal architectures. PHM gear
monitoring work was already in progress when a renewed audit of
the Intel Welding dataset found that some data sessions crossed
the provided partitions and that some proposed test recordings
had already been examined during development. This meant an
independent confirmatory evaluation was not possible on that
assessment population.

The empirical work then shifted to related industrial and robotic
monitoring tasks, while keeping the original architectural research
questions the same.

\ev{Researcher's account of the scope agreed with the supervisor}
{\EvidenceRef{E003}{sections 2.3 and 3.1}}

\ev{Sequence of the gear study and renewed welding audit}
{\EvidenceRef{E004}{sections 11--14}}

PHM2026 focused on building condition targets and sensor-based
prediction. REASSEMBLE was the main experimental setting for
comparing visual and sensor data, fusion methods, and how systems
handle degraded inputs. MSM, ViFailback, and HANDOVER added studies
on milling, robotic manipulation, and human--robot interaction.
More details on dataset feasibility and protocols are in
Appendix~\ref{app:chronology}.

The predictive experiments focus on visual and machine sensor
data. Audio remains relevant to the original research goals and
literature review, but acoustic prediction was not completed.
In REASSEMBLE, the correspondence between action annotations and
acoustic data remained unresolved, and spoken outcomes could have
revealed the prediction labels, introducing a risk of label leakage.
An acoustic prediction branch was therefore not included.

\ev{REASSEMBLE audio feasibility}
{\EvidenceRef{E004}{section 15}}

In this thesis, process-aware fusion means adjusting how predictions
are combined based on the process context and available quality
information. Although the registered title mentions
\textit{physics-informed}, the architectures do not use explicit
differential equations or conservation law constraints.

The empirical work focuses on estimating conditions and classifying
outcomes within set observation periods. Classifying completed
actions and manipulation episodes does not, by itself, provide
prospective anomaly detection or early warning.

Robustness evaluation is limited to experiments where inputs
were removed or degraded in a controlled way. Studies with
external datasets focus on re-using architectures, fitting to
specific datasets, and selected transfer methods, rather than
testing if an unchanged model works in a new domain.

Therefore, the conclusions only apply to the architectures,
prediction targets, and evaluation conditions studied. Independent
welding validation, predictive audio integration, prospective
monitoring, and further operational validation are beyond the
completed empirical scope of this work.


% ============================================================
% 1.6 STRUCTURE OF THE THESIS
% ============================================================

\section{Structure of the Thesis}

Chapter~\ref{ch:literature} is the literature review, setting out the scientific
basis, existing architectural approaches, and remaining research
questions in the area of multimodal industrial monitoring.

Chapter~3 covers the technical underpinnings of the investigated
representation models, the prediction components, and the fusion
mechanisms.

Chapter~\ref{ch:methodology} opens with the five research requirements and maps
them to the empirical studies in Section~\ref{sec:ch4-design}. It then covers
dataset qualification, experimental design, and evaluation procedures.

Chapter~5 looks at condition-target construction and prediction
using PHM2026. Chapter~6 is the main REASSEMBLE chapter, covering
representation comparisons, multimodal fusion, robustness
experiments, and controlled architectural follow-up investigations.

Chapters~7 and~8 examine the application of the investigated
approaches to MSM milling and ViFailback robotic manipulation.
Appendix~\ref{app:handover} provides the supplementary HANDOVER
evaluation.

Chapter~\ref{ch:discussion} brings the empirical findings together
with the relevant literature, answers the research sub-questions,
and looks at the scientific implications across the datasets
considered.

Chapter~\ref{ch:conclusion} covers the supported contributions,
the remaining research questions and the priorities for future work.

The other appendices contain the literature appraisal, dataset
documentation, implementation details, experimental analysis,
research timeline, and reproducibility records.
